\section{Thực nghiệm và kết quả SOTA}
\label{sec::thuc-nghiem}

\subsection*{Thiết lập thực nghiệm}
Để kiểm chứng toàn diện cả bốn tiêu chí của chuẩn \textbf{ISO/IEC 24745} (tính khả hủy, tính không thể liên kết, tính không thể đảo ngược và bảo toàn hiệu năng) cùng khả năng ẩn danh và độ chân thực thị giác, bài báo đánh giá FaceAnonyMixer (tối ưu bằng Adam với $\eta = 0.01$ trong $T = 50$ bước, $N = 5$ biến thể; sử dụng StyleGAN2-FFHQ, e4e, ArcFace ResNet50 và FaRL) trên hai bộ dữ liệu chuẩn là \textbf{CelebA-HQ} và \textbf{VGGFace2}. Quá trình đánh giá được thực hiện trên bốn mô hình nhận dạng độc lập (\textbf{IRSE50}, \textbf{IR152}, \textbf{FaceNet}, \textbf{MobileFace}) tại ba ngưỡng bảo mật $\text{False Matching Rate} \in \{0.1\%, 0.01\%, 0.001\%\}$, đồng thời sử dụng dịch vụ nhận dạng thương mại \textbf{Face++ API} để đo điểm tin cậy tương đồng giữa ảnh bảo vệ và ảnh gốc trong điều kiện ứng dụng thực tế, và đối chiếu trực tiếp với phương pháp học sâu tiên tiến trước đó là \textbf{CanFG}.

\subsection*{Đánh giá tính ẩn danh và tính không thể liên kết}

\begin{table}[H]
\centering
\scriptsize
\setlength{\tabcolsep}{3.5pt}
\caption{Tỷ lệ bảo vệ thành công PSR (\%) đánh giá tính ẩn danh trên CelebA-HQ và VGGFace2 ($\uparrow$ tốt hơn). Nguồn: Bảng 1, bài báo gốc FaceAnonyMixer.}
\label{tab:anonymity_psr}
\begin{tabular}{|l|c|cccc|cccc|c|}
\hline
\multirow{2}{*}{\textbf{Phương pháp}} & \multirow{2}{*}{\textbf{FMR}} & \multicolumn{4}{c|}{\textbf{CelebA-HQ}} & \multicolumn{4}{c|}{\textbf{VGGFace2}} & \multirow{2}{*}{\textbf{Trung bình}} \\
\cline{3-10}
 & & IRSE50 & IR152 & FaceNet & MobileFace & IRSE50 & IR152 & FaceNet & MobileFace & \\
\hline
\multirow{3}{*}{CanFG} & 0.1\% & 71.85 & 70.21 & 81.48 & 34.81 & 1.09 & 15.27 & 30.68 & 0.07 & 38.18 \\
 & 0.01\% & 60.49 & 54.40 & 71.44 & 22.88 & 0.36 & 8.54 & 21.27 & 0.05 & 29.93 \\
 & 0.001\% & 50.37 & 40.00 & 60.58 & 14.65 & 0.14 & 4.20 & 15.85 & 0.01 & 23.22 \\
\hline
\multirow{3}{*}{\textbf{FaceAnonyMixer}} & 0.1\% & \textbf{78.68} & \textbf{95.47} & \textbf{99.51} & \textbf{46.83} & \textbf{98.61} & \textbf{99.07} & \textbf{98.15} & \textbf{93.52} & \textbf{88.73} \\
 & 0.01\% & \textbf{65.58} & \textbf{84.28} & \textbf{97.28} & \textbf{29.14} & \textbf{96.76} & \textbf{97.69} & \textbf{95.37} & \textbf{85.19} & \textbf{81.41} \\
 & 0.001\% & \textbf{61.62} & \textbf{65.68} & \textbf{91.60} & \textbf{24.85} & \textbf{90.67} & \textbf{92.06} & \textbf{88.81} & \textbf{74.09} & \textbf{73.67} \\
\hline
\end{tabular}
\end{table}

\begin{table}[H]
\centering
\scriptsize
\setlength{\tabcolsep}{3.5pt}
\caption{Tỷ lệ bảo vệ thành công PSR (\%) đa khóa đánh giá tính không thể liên kết ($\uparrow$ tốt hơn). Nguồn: Bảng 2, bài báo gốc FaceAnonyMixer.}
\label{tab:unlinkability_psr}
\begin{tabular}{|l|c|cccc|cccc|c|}
\hline
\multirow{2}{*}{\textbf{Phương pháp}} & \multirow{2}{*}{\textbf{FMR}} & \multicolumn{4}{c|}{\textbf{CelebA-HQ}} & \multicolumn{4}{c|}{\textbf{VGGFace2}} & \multirow{2}{*}{\textbf{Trung bình}} \\
\cline{3-10}
 & & IRSE50 & IR152 & FaceNet & MobileFace & IRSE50 & IR152 & FaceNet & MobileFace & \\
\hline
\multirow{3}{*}{CanFG} & 0.1\% & 53.50 & 56.13 & 72.51 & 22.80 & 4.36 & 8.83 & 25.18 & 2.41 & 30.72 \\
 & 0.01\% & 40.49 & 39.59 & 58.27 & 11.28 & 2.41 & 4.63 & 16.71 & 1.13 & 21.81 \\
 & 0.001\% & 29.05 & 24.44 & 45.10 & 5.76 & 0.19 & 2.39 & 11.8 & 0.21 & 14.87 \\
\hline
\multirow{3}{*}{\textbf{FaceAnonyMixer}} & 0.1\% & \textbf{86.21} & \textbf{95.56} & \textbf{97.58} & \textbf{47.10} & \textbf{98.35} & \textbf{98.75} & \textbf{97.89} & \textbf{92.20} & \textbf{89.21} \\
 & 0.01\% & \textbf{66.45} & \textbf{77.82} & \textbf{95.56} & \textbf{30.97} & \textbf{95.15} & \textbf{96.14} & \textbf{95.87} & \textbf{85.21} & \textbf{80.40} \\
 & 0.001\% & \textbf{45.08} & \textbf{60.48} & \textbf{91.53} & \textbf{19.27} & \textbf{88.17} & \textbf{90.47} & \textbf{82.69} & \textbf{77.34} & \textbf{69.38} \\
\hline
\end{tabular}
\end{table}

Từ Bảng~\ref{tab:anonymity_psr} và Bảng~\ref{tab:unlinkability_psr}, FaceAnonyMixer vượt trội hoàn toàn CanFG ở mọi mô hình và mọi ngưỡng bảo mật. Cụ thể, khi đặt ngưỡng nhận diện ở mức chỉ cho phép nhận nhầm người lạ $0.1\%$, tỷ lệ hệ thống không nhận ra ảnh bảo vệ là người gốc tăng từ $38.18\%$ lên $88.73\%$ ($+50.55\%$); đồng thời tỷ lệ hai ảnh bảo vệ của cùng một người nhưng khác khóa không bị nhận ra là cùng một người tăng từ $30.72\%$ lên $89.21\%$ ($+58.49\%$). Trên tập dữ liệu phức tạp VGGFace2, CanFG gần như thất bại (tỷ lệ che giấu được danh tính gốc chỉ còn $0.01\%\text{--}30.68\%$), trong khi FaceAnonyMixer vẫn đạt trên $90\%$ ở hầu hết các mô hình. Sự khác biệt này bắt nguồn trực tiếp từ hai thay đổi kiến trúc cốt lõi: (1) CanFG biến đổi trực tiếp trên mạng U-Net có liên kết tắt (\textit{skip-connections}) giữ lại nhiều chi tiết không gian của ảnh gốc và thiếu hàm tối ưu trực giao, nên danh tính cũ vẫn bị rò rỉ và các khóa khác nhau cho ra khuôn mặt na ná nhau; (2) ngược lại, FaceAnonyMixer \textbf{thay thế triệt để các tầng danh tính ($3\text{--}7$)} trong không gian $\mathcal{W}^+$ giàu biểu diễn của StyleGAN2 kết hợp hàm $\mathcal{L}_{\text{anon}}$ ép vector đặc trưng tiến tới trực giao với ảnh gốc ($m=0$). Nhờ đó, mỗi khóa $\mathbf{k}$ đưa khuôn mặt sang một danh tính giả độc lập.

\subsection*{Đánh giá thực tế, chống đảo ngược, hiệu năng nhận dạng và độ ổn định khóa}

\begin{table}[H]
\centering
\scriptsize
\setlength{\tabcolsep}{2pt}
\begin{minipage}[t]{0.45\textwidth}
\centering
\caption{Điểm tin cậy trên Face++ API ($\downarrow$ tốt hơn). Nguồn: Bảng 3, bài báo gốc.}
\label{tab:facepp}
\begin{tabular}{|l|ccc|}
\hline
\textbf{Phương pháp} & \textbf{CelebA-HQ} & \textbf{VGGFace2} & \textbf{TB} \\
\hline
CanFG & 71.37 & 70.78 & 71.07 \\
\textbf{FaceAnonyMixer} & \textbf{53.73} & \textbf{31.51} & \textbf{58.37} \\
\hline
\end{tabular}
\end{minipage}
\hfill
\begin{minipage}[t]{0.52\textwidth}
\centering
\caption{Hiệu năng xác thực EER ($\downarrow$ tốt hơn) và AUC ($\uparrow$ tốt hơn). Nguồn: Bảng 5, bài báo gốc.}
\label{tab:utility}
\begin{tabular}{|l|cc|cc|}
\hline
\multirow{2}{*}{\textbf{Phương pháp}} & \multicolumn{2}{c|}{\textbf{CelebA-HQ}} & \multicolumn{2}{c|}{\textbf{VGGFace2}} \\
\cline{2-5}
 & EER$\downarrow$ & AUC$\uparrow$ & EER$\downarrow$ & AUC$\uparrow$ \\
\hline
Original (Gốc) & 0.027 & 0.990 & 0.074 & 0.964 \\
CanFG & 0.045 & 0.988 & 0.101 & 0.951 \\
\textbf{FaceAnonyMixer} & \textbf{0.019} & \textbf{0.997} & \textbf{0.059} & \textbf{0.975} \\
\hline
\end{tabular}
\end{minipage}
\end{table}

Các kết quả bổ trợ dưới đây củng cố tính toàn diện của phương pháp:
\begin{itemize}
    \item \textbf{Kiểm tra tính ứng dụng bằng API nhận dạng đời thực (Bảng~\ref{tab:facepp}):} Điểm tin cậy tương đồng trên Face++ giảm từ $71.07$ (CanFG) xuống $58.37$ (riêng VGGFace2 giảm sâu từ $70.78$ còn $31.51$), chứng tỏ việc tối ưu trong $\mathcal{W}^+$ đánh lừa hiệu quả ngay cả các hệ thống thương mại không rõ kiến trúc.
    \item \textbf{Chống đảo ngược ngay cả khi lộ khóa (Bảng~\ref{tab:noninvertibility}, Phụ lục~\ref{app:extra_tables}):} Ngay cả khi kẻ tấn công biết cả khóa $\mathbf{k}$ lẫn ảnh bảo vệ để thay lại tầng giữa hoặc huấn luyện mạng dịch ảnh ngược \textit{pix2pix}, tỷ lệ ảnh tái tạo không bị nhận ra là người gốc vẫn đạt $93.80\%\text{--}99.26\%$ tại ngưỡng chỉ cho phép nhận nhầm người lạ $0.1\%$. Lý do là phần mã định danh gốc $\mathbf{z}_r^{(3:7)}$ đã bị loại bỏ hoàn toàn trước khi tối ưu phi tuyến qua 50 bước lặp, khiến ánh xạ ngược trở nên bất khả thi.
    \item \textbf{Bảo toàn nhận dạng và chất lượng ảnh (Bảng~\ref{tab:utility}):} CanFG làm suy giảm khả năng xác thực so với ảnh gốc (EER trên CelebA-HQ tăng từ $0.027 \rightarrow 0.045$; trên VGGFace2 tăng từ $0.074 \rightarrow 0.101$) vì tối ưu đối nghịch trên U-Net không có ràng buộc giữ biểu cảm, hướng nhìn gây méo đặc trưng (FID $= 82.52$). Ngược lại, FaceAnonyMixer giảm EER xuống \textbf{0.019} (CelebA-HQ) và \textbf{0.059} (VGGFace2) --- thậm chí tốt hơn cả ảnh gốc chưa bảo vệ --- nhờ $\mathcal{L}_{\text{idp}}$ gom chặt các ảnh bảo vệ cùng người dưới cùng khóa trong không gian ArcFace, trong khi $\mathcal{L}_{\text{attr}}$ (FaRL) và bộ sinh StyleGAN2 giữ biểu cảm tự nhiên với FID chỉ còn \textbf{37.73}.
    \item \textbf{Ổn định trước sự thay đổi khóa (Bảng~\ref{tab:key_robustness}, Phụ lục~\ref{app:extra_tables}):} Qua $5$ khóa ngẫu nhiên khác nhau, tỷ lệ bảo vệ thành công trên IR152 duy trì trung bình $93.51\%$ với độ lệch chuẩn rất nhỏ ($\text{Std.} = 2.33$), khẳng định hiệu quả bảo vệ không phụ thuộc vào việc chọn khóa đặc biệt.
\end{itemize}